%!TEX root = ../dokumentation.tex
%
% Alle Stil-Einstellungen der Arbeit, in Abschnitten:
%    1 Seite und Zeilen          6 Quellcode
%    2 Schrift                   7 Sonstiges (Fußnoten, Entwurfsnotizen)
%    3 Sprache                   8 Links und PDF-Infos (hyperref)
%    4 Farben                    9 Zitierweise
%    5 Abbildungen, Tabellen,   10 Glossar und Abkürzungen
%      Diagramme (TikZ-Stile)
% Persönliche Angaben stehen nicht hier, sondern oben in dokumentation.tex.
% Reihenfolge beachten: hyperref (8) lädt spät; 9 und 10 brauchen es.
%

%% ---- 1 Seite und Zeilen --------------------------------------------
% Seitenränder: fullpage = rundum ca. 2,5 cm. Feste Maße stattdessen:
%   \usepackage[left=3cm,right=2.5cm,top=2.5cm,bottom=2.5cm]{geometry}
\usepackage{fullpage}

% Zeilenabstand: onehalfspacing = 1,5-zeilig; singlespacing, doublespacing
% oder freier Wert mit \setstretch{1.3} nach dem Laden.
\usepackage[onehalfspacing]{setspace}

% Hurenkinder und Schusterjungen verhindern
% http://projekte.dante.de/DanteFAQ/Silbentrennung
\clubpenalty=10000
\widowpenalty=10000
\displaywidowpenalty=10000

%% ---- 2 Schrift -----------------------------------------------------
\usepackage[utf8]{inputenc}
\usepackage[T1]{fontenc}

% Schriftart Palatino. Alternativen (eine Zeile statt palatino):
%   \usepackage{lmodern}               Latin Modern (LaTeX-Standard)
%   \usepackage{libertine}             Linux Libertine
%   \usepackage{mathptmx}              Times
%   \usepackage{newpxtext,newpxmath}   Palatino-Klon mit passender Mathe
\usepackage{palatino}

% Zeilenumbruch und mehr
\usepackage[activate]{microtype}

% Überschriften setzt KOMA-Script serifenlos fett. In der Textschrift:
%\addtokomafont{disposition}{\rmfamily}

%% ---- 3 Sprache -----------------------------------------------------
% Die Arbeit hat genau eine Sprache: ./set-language.sh de|en setzt
% \englischtrue / \englischfalse in dokumentation.tex und das passende Gerüst.
% Deckblatt, Erklärung, Verzeichnisse, Zitate und Verweise folgen dem Schalter.
% \deen{deutsch}{english}: Text je nach Sprache, für Vorlagenteile wie das Deckblatt
\newcommand{\deen}[2]{\ifenglisch#2\else#1\fi}
% main= explizit: 'ngerman' steht auch in den Klassenoptionen und wird
% global durchgereicht, wodurch die reine Reihenfolge nicht mehr entscheidet.
\ifenglisch
	\usepackage[ngerman,main=english]{babel}
\else
	\usepackage[english,main=ngerman]{babel}
\fi

% Anführungszeichen: \enquote{...} setzt sie passend zur Sprache
\usepackage[babel,german=quotes]{csquotes}

% Überschriften, die babel nicht kennt
\ifenglisch
	\newcommand{\abkuerzungstitel}{List of Abbreviations}
	\newcommand{\glossartitel}{Glossary}
\else
	\newcommand{\abkuerzungstitel}{Abkürzungsverzeichnis}
	\newcommand{\glossartitel}{Glossar}
\fi

% Index-Erstellung
\usepackage{makeidx}

%% ---- 4 Farben ------------------------------------------------------
% xcolor, not color: listings and TikZ load xcolor later, and definitions made
% with plain `color` then raise "Incompatible color definition".
\usepackage{xcolor}
% Farbe aller Links im PDF. Für den Druck schwarz: {rgb}{0,0,0}
\definecolor{LinkColor}{rgb}{0,0,0.2}
% Grundfarbe aller TikZ-Diagramme (Abschnitt 5), z.B. blue!60!black
\colorlet{akzent}{black}

%% ---- 5 Abbildungen, Tabellen, Diagramme ----------------------------
% Grafiken liegen in images/; Dateiname ohne Endung angeben.
\usepackage{graphicx}
\graphicspath{{images/}}

% Gleitumgebungen (Bilder, Tabellen, usw\ldots) lassen sich mit H an genau der
% definierten Stelle platzieren
\usepackage{float}

% für textumflossene Grafiken
\usepackage{wrapfig}

% Tabellen: booktabs (\toprule, \midrule, \bottomrule), tabularx (Spalten
% füllen die Breite), array (vertikale Platzierung), longtable (über Seiten).
\usepackage{booktabs}
\usepackage{tabularx}
\usepackage{array}
\usepackage{longtable}
% Gilt für ALLE Tabellen, auch das Deckblatt: Spaltenabstand und Zeilenhöhe
\setlength{\tabcolsep}{10pt}
\renewcommand{\arraystretch}{1.5}

% Formeln
\usepackage{amsmath}
%\usepackage{amssymb}

% Pakete um Textteile drehen zu können, oder eine Seite Querformat anzeigen kann.
%\usepackage{rotating}
%\usepackage{lscape}

% TikZ-Diagramme mit einem gemeinsamen Stilbaukasten, damit alle Diagramme
% gleich aussehen. Farbe zentral über 'akzent' (Abschnitt 4).
%   flow   = Grundstil fürs Bild:  \begin{tikzpicture}[flow]
%   fio    = Ein-/Ausgabe (gerundet, leicht gefüllt)
%   fstage = Verarbeitungsschritt
%   fkey   = hervorgehobener Schritt (dunkler gefüllt)
%   fmuted = Nebenschritt (heller Rahmen, graue Schrift)
%   fgate  = Entscheidung (Raute)
%   fedge  = Pfeil;  fcond = bedingter Pfeil (gestrichelt)
%   flab   = Beschriftung an einem Pfeil;  fnote = kleine Randnotiz
% Beispiel: content/beispiele.tex
\usepackage{tikz}
\usetikzlibrary{arrows.meta,positioning,shapes.geometric,calc}
% Diagramme nie breiter als der Satzspiegel: max width verkleinert nur,
% vergroessert aber nie (im Gegensatz zu \resizebox).
\usepackage{adjustbox}
\tikzset{
  flow/.style={font=\footnotesize\sffamily, line cap=round},
  fio/.style={draw=akzent!70, fill=akzent!8, rounded corners=2.2mm, inner xsep=2.4mm, inner ysep=1.5mm, align=center},
  fstage/.style={draw=akzent, line width=0.7pt, fill=white, rounded corners=0.6mm, inner xsep=2.4mm, inner ysep=1.6mm, align=center},
  fkey/.style={fstage, fill=akzent!15},
  fmuted/.style={draw=akzent!35, fill=white, rounded corners=0.6mm, inner xsep=2.2mm, inner ysep=1.5mm, align=center, text=black!60},
  fgate/.style={draw=akzent!50, fill=akzent!4, diamond, aspect=2.4, inner sep=0.8mm, align=center},
  fedge/.style={-{Stealth[length=2mm]}, draw=akzent!55, semithick},
  fcond/.style={fedge, dashed},
  flab/.style={font=\scriptsize\sffamily, text=black!60, inner sep=1pt, fill=white},
  fnote/.style={font=\scriptsize\sffamily, text=black!45, inner sep=1pt},
}

%% ---- 6 Quellcode ---------------------------------------------------
% listings statt minted: braucht kein -shell-escape und kein Python.
% Sprache pro Listing ändern: \begin{lstlisting}[language=Java]
\usepackage{listings}
\definecolor{noteframe}{gray}{0.55}
\lstset{
  basicstyle=\ttfamily\scriptsize,
  keywordstyle=\color{black}\bfseries,
  commentstyle=\itshape\color{noteframe},
  stringstyle=\color{black},
  numbers=left, numberstyle=\tiny\color{noteframe}, numbersep=7pt,
  frame=leftline, framesep=6pt, rulecolor=\color{noteframe},
  breaklines=true, breakatwhitespace=false, breakindent=12pt,
  showstringspaces=false, tabsize=2, columns=fullflexible,
  language=Python, upquote=true, aboveskip=10pt, belowskip=10pt,
  captionpos=b, xleftmargin=10pt,
  % listings reads its body byte by byte, so every non-ASCII character that
  % appears in the source code needs an explicit mapping.
  extendedchars=true, inputencoding=utf8,
  literate={ä}{{\"a}}1 {ö}{{\"o}}1 {ü}{{\"u}}1
           {Ä}{{\"A}}1 {Ö}{{\"O}}1 {Ü}{{\"U}}1
           {ß}{{\ss}}1 {–}{{--}}1 {—}{{---}}1
           {─}{{-}}1,
}
\renewcommand{\lstlistingname}{Listing}
\ifenglisch\else\renewcommand{\lstlistlistingname}{Quellcodeverzeichnis}\fi

%% ---- 7 Sonstiges ---------------------------------------------------
% Fussnoten
\usepackage[perpage, hang, multiple, stable]{footmisc}

% für die Darstellung des Euro-Symbols: \euro{}
\usepackage[right]{eurosym}

% Entwurfsnotizen: \begin{notiz} ... \end{notiz} wird rot gedruckt.
% Vor der Abgabe \excludecomment{notiz} einkommentieren: alle Notizen verschwinden.
% (Nicht 'k' nennen: \k ist der Ogonek-Akzent für ą/ę, den hyperref neu definiert.)
\usepackage{comment}
% Zwischendatei als dokumentation.cut statt comment.cut, damit latexmk -c sie aufräumt
\def\DefaultCutFileName{\def\CommentCutFile{\jobname.cut}}
\specialcomment{notiz}{\begingroup\color{red}}{\endgroup}
%\excludecomment{notiz}

%% ---- 8 Links und PDF-Infos -----------------------------------------
\usepackage[%
	pdftitle={\titel},
	pdfauthor={\autor},
	pdfsubject={\arbeit},
	pdfcreator={pdflatex, LaTeX with KOMA-Script},
	pdfpagemode=UseOutlines, % Beim Oeffnen Inhaltsverzeichnis anzeigen
	pdfdisplaydoctitle=true, % Dokumenttitel statt Dateiname anzeigen.
]{hyperref}

% (Farb-)einstellungen für die Links im PDF.
% Ganz ohne farbige Links: colorlinks=false, hidelinks
\hypersetup{%
	colorlinks=true, % Aktivieren von farbigen Links im Dokument
	linkcolor=LinkColor, % Farbe festlegen
	citecolor=LinkColor,
	filecolor=LinkColor,
	menucolor=LinkColor,
	urlcolor=LinkColor,
	bookmarksnumbered=true % Überschriftsnummerierung im PDF Inhalt anzeigen.
}
\ifenglisch\hypersetup{pdflang=en}\else\hypersetup{pdflang=de}\fi

% \autoref{lst:...} schreibt "Listing" vor die Nummer
\providecommand*{\lstlistingautorefname}{Listing}
% Englisch: hyperref schreibt "chapter 2"; großgeschrieben passt es auch am Satzanfang
\ifenglisch
	\AtBeginDocument{%
		\def\chapterautorefname{Chapter}%
		\def\sectionautorefname{Section}%
		\def\subsectionautorefname{Section}%
		\def\subsubsectionautorefname{Section}%
		\def\appendixautorefname{Appendix}%
	}
\fi

%nur ein latex-Durchlauf für die Aktualisierung von Verzeichnissen nötig
\usepackage{bookmark}

%% ---- 9 Zitierweise -------------------------------------------------
% Literaturverweise nach Harvard (Autor, Jahr) mit agsm.bst
\usepackage[dcucite]{harvard}
\ifenglisch\renewcommand{\harvardand}{and}\else\renewcommand{\harvardand}{und}\fi
% agsm.bst emits \harvardurl{<raw url>}; the original grabs its argument
% normally, so any '_' in a URL raises "Missing $ inserted".
% Handing the braces straight to \url makes it read them verbatim.
\AtBeginDocument{\renewcommand{\harvardurl}{\textbf{URL:} \url}}
% Ab drei Autoren immer (Autor et al., Jahr)
\citationmode{abbr}

% Zitierbefehle (natbib-Namen; harvard.sty kennt dieselben Formen unter
% anderem Namen):
%   \citep{key}          -> (Autor, Jahr)        == \cite{key}
%   \citep[S.~12]{key}   -> (Autor, Jahr, S. 12)
%   \citet{key}          -> Autor (Jahr)         == \citeasnoun{key}
\newcommand{\citep}[2][]{\if\relax\detokenize{#1}\relax\cite{#2}\else\cite[#1]{#2}\fi}
\newcommand{\citet}[2][]{\if\relax\detokenize{#1}\relax\citeasnoun{#2}\else\citeasnoun[#1]{#2}\fi}

%% ---- 10 Glossar und Abkürzungen ------------------------------------
% Glossar
\usepackage[
	nonumberlist, %keine Seitenzahlen anzeigen
	%acronym,      %ein Abkürzungsverzeichnis erstellen
	%section,     %im Inhaltsverzeichnis auf section-Ebene erscheinen
	toc,          %Einträge im Inhaltsverzeichnis
]{glossaries}

% Akronyme: \ac{API} -- beim ersten Mal ausgeschrieben (als Fußnote)
\usepackage[printonlyused,footnote]{acronym}
